Adaptive financial systems can affect the environment from which they learn, but deployment traces and standalone model scores do not by themselves establish market-level effects. We present an evidence-bounded study combining a deployed two-role decision-system audit, a matched-input response analysis, and controlled market simulations. The operational audit covers 28,165 events and 229 entry-to-execution-response pairs, while identifying missing portfolio-model responses that prevent complete two-role replay. Across 5,998 common exported inputs, four model labels exhibit pairwise directional agreement of 52.2--80.7 percent; these measurements describe the export rather than original matched inference or trading accuracy. We implement a single-asset market with explicit inventory, cash, trading costs, order limits, and price feedback, and run 1,312 episodes spanning capital, depth, population size, trade limits, behavioral perturbations, and external shocks. In a 20-seed paired experiment, raising adaptive capital participation from 25 to 75 percent changes the incremental shock effect on maximum drawdown in different directions across assumed response profiles. A selective threshold-and-drawdown profile does not uniformly reduce shock amplification. Accounting checks and executable artifacts make the simulation findings reproducible under their stated assumptions. The results distinguish observed response diversity from assumed policy diversity and show why claims about adaptive participation require explicit counterfactual mechanisms. They do not establish risk-adjusted alpha, endogenous liquidity withdrawal, or a calibrated market-crash boundary.